\chapter{What a Signal Is}

\chapterline

Before we touch a single Fourier integral, before we sample anything or design a filter, we
need to agree on what we are even looking at. The whole edifice of signal processing rests on
one deceptively simple object: the \emph{signal}. Get a clear, geometric picture of what a
signal is, and everything that follows -- frequency, the Fourier transform, filtering,
sampling -- becomes a story about \emph{shapes} and \emph{measurements} rather than a wall of
formulas. Get a muddy picture, and you will spend the rest of the book pushing symbols around
without ever seeing what they mean.

This chapter builds that picture from the ground up. We will see that a signal is, at heart,
just a function: a rule that assigns a value to each instant of time (or each point in space).
We will look at real signals you can hear, see, and wear on your wrist. We will draw the great
dividing lines of the field -- one dimension versus two, single channel versus many, analog
versus digital -- and we will state the five central questions that the rest of the book exists
to answer. And throughout, we will keep one concrete project in view: a wrist- or forearm-worn
\textbf{gesture controller} that turns the motion of your hand into music.

\section{A signal is a function carrying information}
\label{ch01:defn}

Pluck a guitar string. The string vibrates; it shoves the air; the air pressure at your eardrum
rises and falls hundreds of times a second. If you could write down the air pressure at every
instant, you would have a number for each moment in time -- a \emph{function of time}. That
function \emph{is} the sound. Nothing is lost in describing the sound this way; the wiggling
number contains everything your ear receives.

\begin{definition}[Signal, working definition]
\label{ch01:signal-defn}
A \textbf{signal} is a function $x$ that maps an independent variable (such as time $t$ or
position) to a value $x(t)$, where the variation of that value carries \emph{information}.
\end{definition}

Two words in that definition deserve a second look. The first is \emph{function}. A signal is
not the wire it travels on, not the sensor that produced it, not the file it is stored in. It is
the mathematical object: an assignment of a value to each point of some domain. When we write
\[
  x : \R \to \R, \qquad t \mapsto x(t),
\]
we mean exactly the air-pressure-versus-time idea, abstracted so we can do mathematics to it.

The second word is \emph{information}. A constant function $x(t) = 5$ for all $t$ is a perfectly
good function, but as a signal it is dull: it never changes, so it tells you nothing beyond its
one value. Information lives in \emph{variation} -- in how the signal differs from instant to
instant. This is why the entire field is obsessed with \emph{change}, and why frequency (the
rate of a particular kind of change) will turn out to be the master concept of the book.

\begin{intuitionbox}
Think of a signal as a recording of "what was the value, moment by moment." A thermometer you
read once an hour produces a signal: temperature versus time. A microphone produces a signal:
pressure versus time. A photograph produces a signal: brightness versus position. In every case
the same skeleton appears -- an independent variable (the input) and a dependent value (the
output) that wiggles in a way that means something.
\end{intuitionbox}

\subsection*{The graph is the signal}
The most important habit you can build right now is to \emph{see the graph}. Whenever someone
says "a signal," picture its plot: the horizontal axis is the independent variable, the vertical
axis is the value. Figure~\ref{ch01:fig-tone} shows the simplest interesting signal -- a pure
tone, a single sinusoid. Stare at it. The height at each $t$ is the value $x(t)$; the wiggling
is the information; the rate of wiggling is the frequency.

\begin{figure}[H]
\centering
\begin{tikzpicture}
\begin{axis}[
  width=12cm, height=5cm,
  xlabel={time $t$ (seconds)}, ylabel={$x(t)$},
  xmin=0, xmax=1, ymin=-1.4, ymax=1.4,
  axis lines=left,
  domain=0:1, samples=200,
  every axis plot/.append style={thick},
  enlargelimits=false,
]
\addplot[sigblue] {sin(deg(2*pi*4*x))};
\end{axis}
\end{tikzpicture}
\caption{A pure tone: $x(t)=\sin(2\pi f t)$ with $f=4$~Hz. The value at each instant is the
height of the curve; the four full up-and-down cycles per second \emph{are} its frequency.}
\label{ch01:fig-tone}
\end{figure}

\begin{whybox}
Why insist on the graph instead of the formula $x(t)=\sin(2\pi f t)$? Because most real signals
have no tidy formula at all -- speech, an accelerometer trace during a wrist flick, a
heartbeat. They exist only as data: a long list of measured values. The graph is the one
representation that works for \emph{every} signal, formula or not. Geometry is the common
language; the algebra is a convenience that happens to apply to special cases.
\end{whybox}

\section{A gallery of real signals}
\label{ch01:gallery}

Definitions are cheap; intuition comes from examples. Here is a small gallery of signals you
will meet again and again in this book. Notice that they are wildly different physically -- some
you hear, some you see, some you feel -- yet they are all the same kind of mathematical object.

\paragraph{Audio.} Sound is variation in air pressure over time. A microphone converts that
pressure into a voltage; the signal is voltage (proportional to pressure) versus time. Human
hearing runs from roughly 20~Hz to 20~kHz, which already tells you something quantitative: the
fastest meaningful wiggles in music happen tens of thousands of times per second.

\paragraph{Images.} A grayscale photograph is brightness as a function of \emph{position}:
$x(u,v)$ where $(u,v)$ names a pixel. Here the independent variable is two-dimensional space, not
time. The "information" is the spatial pattern -- edges, textures, shapes.

\paragraph{Accelerometer traces.} Strap a tiny accelerometer to your wrist and it reports
acceleration along three axes, dozens or hundreds of times per second. Wave your hand and you
get three wiggling curves. This is exactly the raw material of our gesture controller, so we
will return to it constantly.

\paragraph{Biosignals.} Your body is a noisy broadcaster. The electrocardiogram (ECG) is the
heart's electrical activity versus time. The electromyogram (EMG) is the electrical activity of
muscles -- and, crucially for finger-level sensing, the forearm muscles that move your fingers
sit right under a wristband. We will lean hard on this fact in Chapter 37.

Figure~\ref{ch01:fig-gallery} sketches three of these side by side so you can feel both their
diversity and their shared skeleton.

\begin{figure}[H]
\centering
\begin{tikzpicture}
% --- panel 1: audio-like ---
\begin{axis}[
  name=ax1, width=5.6cm, height=3.4cm,
  title={(a) audio (pressure)},
  xlabel={$t$}, xtick=\empty, ytick=\empty,
  xmin=0, xmax=1, ymin=-1.3, ymax=1.3,
  axis lines=left, domain=0:1, samples=180,
]
\addplot[sigblue, thick] {0.6*sin(deg(2*pi*9*x)) + 0.4*sin(deg(2*pi*23*x))};
\end{axis}
% --- panel 2: accelerometer-like ---
\begin{axis}[
  name=ax2, at={($(ax1.east)+(0.9cm,0)$)}, anchor=west,
  width=5.6cm, height=3.4cm,
  title={(b) wrist accel (1 axis)},
  xlabel={$t$}, xtick=\empty, ytick=\empty,
  xmin=0, xmax=1, ymin=-1.3, ymax=1.3,
  axis lines=left, domain=0:1, samples=180,
]
\addplot[targetgreen, thick] {exp(-6*(x-0.45)^2)*sin(deg(2*pi*5*x)) };
\end{axis}
% --- panel 3: ECG-like ---
\begin{axis}[
  name=ax3, at={($(ax2.east)+(0.9cm,0)$)}, anchor=west,
  width=5.6cm, height=3.4cm,
  title={(c) ECG (heart)},
  xlabel={$t$}, xtick=\empty, ytick=\empty,
  xmin=0, xmax=1, ymin=-0.6, ymax=1.3,
  axis lines=left, domain=0:1, samples=200,
]
\addplot[specred, thick] {1.1*exp(-700*(x-0.25)^2) + 1.1*exp(-700*(x-0.75)^2) - 0.18*exp(-300*(x-0.18)^2) - 0.18*exp(-300*(x-0.68)^2)};
\end{axis}
\end{tikzpicture}
\caption{Three signals, one skeleton. (a) A bit of audio (a sum of two tones). (b) One axis of
a wrist accelerometer during a single flick -- a burst that grows and fades. (c) A stylized ECG
with two heartbeats. Different physics, same object: a value versus time.}
\label{ch01:fig-gallery}
\end{figure}

\begin{goalbox}
Our running goal is a \textbf{wrist/forearm gesture controller} that outputs MIDI -- in plain
terms, a band you wear that watches your hand and plays an instrument. Panel (b) above is the
heart of it: every gesture you make becomes a little story written in acceleration (and rotation
and muscle activity) over time. The entire job of the controller is to take those wiggling
curves and decide, in real time, \emph{what did the hand just do?} -- then emit the right
musical note or control change. Every chapter of this book is a tool for reading panel (b)
correctly. Keep that image in mind: we are learning to read the handwriting of motion.
\end{goalbox}

\section{The dimensionality of the domain: 1D, 2D, and beyond}
\label{ch01:dimension}

Signals are sorted first by the shape of their \emph{input}, the independent variable.

\paragraph{One-dimensional (1D).} The input is a single number -- almost always time. Audio,
ECG, a single accelerometer axis: all 1D. The graph is a curve. This is where we will spend most
of the book, because the ideas are cleanest here and generalize cleanly.

\paragraph{Two-dimensional (2D).} The input is a pair, typically a position $(u,v)$. An image is
the canonical 2D signal: $x(u,v)$ is the brightness at pixel $(u,v)$. The "graph" is now a
surface, or more usefully a picture. Everything we learn about 1D frequency will reappear in 2D
as \emph{spatial} frequency -- fine textures are high spatial frequency, broad gradients are low.
Chapters 30--33 are devoted to this.

\paragraph{Higher dimensions.} Video is a 3D signal $x(u,v,t)$: brightness as a function of two
space coordinates and time. A volumetric medical scan is $x(u,v,w)$ over three space
coordinates. The machinery extends, but the intuition you build in 1D and 2D carries you the
rest of the way.

\begin{focusbox}
Do not confuse the dimension of the \emph{domain} (the input) with the dimension of the
\emph{value} (the output). A grayscale image is a 2D-domain, 1D-value signal: two-number input,
one-number brightness out. A color image is 2D-domain, 3D-value (red, green, blue). Time is 1D
no matter how many sensors you attach -- a hundred microphones still listen along the single
axis of time. The next section is about that output side: how many values come out at each
instant.
\end{focusbox}

\section{Single channel versus multichannel}
\label{ch01:channels}

So far each signal has produced \emph{one} value at each instant. But many real systems produce
several at once. A stereo recording gives two values per instant (left ear, right ear). A
three-axis accelerometer gives three. A modern wrist band might stream nine numbers per instant:
three of acceleration, three of rotation rate (gyroscope), three of magnetic field
(magnetometer).

We call each of these parallel streams a \textbf{channel}. A \emph{multichannel} signal is best
thought of as a \emph{vector-valued} function of time:
\[
  \vec{x}(t) = \big(x_1(t),\, x_2(t),\, \dots,\, x_C(t)\big),
\]
where $C$ is the number of channels. At each instant the signal is a point in $\R^C$; as time
runs, that point traces a path. For the nine-axis wrist band, $\vec{x}(t)$ is a moving point in
nine-dimensional space, and a "gesture" is a recognizable shape of that path.

Figure~\ref{ch01:fig-trajectory} shows the idea for a two-channel signal. On the left the two
channels are drawn the usual way, each as its own curve against time. On the right the same data
is drawn as a single path -- channel~1 on the horizontal axis, channel~2 on the vertical -- with
time hidden inside the motion along the curve. A smooth circular wave of the hand produces a
loop; that loop \emph{is} the gesture, seen as a shape in channel space.

\begin{figure}[H]
\centering
\begin{tikzpicture}
\begin{axis}[
  name=axt, width=6.4cm, height=4.2cm,
  title={(a) two channels vs.\ time},
  xlabel={$t$}, ylabel={value}, xtick=\empty,
  ytick={-1,0,1}, yticklabels={$-1$,$0$,$1$},
  xmin=0, xmax=1, ymin=-1.3, ymax=1.3,
  axis lines=left, domain=0:1, samples=120,
  legend style={at={(0.5,-0.18)}, anchor=north, legend columns=2, draw=none,
    font=\footnotesize, /tikz/every even column/.append style={column sep=0.3cm}},
]
\addplot[sigblue, thick] {cos(deg(2*pi*2*x))};
\addlegendentry{$x_1(t)$}
\addplot[targetgreen, thick] {sin(deg(2*pi*2*x))};
\addlegendentry{$x_2(t)$}
\end{axis}
\begin{axis}[
  name=axp, at={($(axt.east)+(1.4cm,0)$)}, anchor=west,
  width=4.2cm, height=4.2cm,
  title={(b) path in channel space},
  xlabel={$x_1$}, ylabel={$x_2$},
  xmin=-1.3, xmax=1.3, ymin=-1.3, ymax=1.3,
  axis lines=middle, domain=0:1, samples=120,
  xtick={-1,1}, ytick={-1,1},
]
\addplot[structurepurple, thick] ({cos(deg(2*pi*2*x))}, {sin(deg(2*pi*2*x))});
\end{axis}
\end{tikzpicture}
\caption{Two views of one two-channel signal. (a) Each channel as its own curve against time.
(b) The same signal as a point $(x_1(t),x_2(t))$ tracing a path in channel space; here a quarter
of the run has gone by before the loop closes. A faster wave runs around the same loop more times
in the same window. Exercise~\ref{ch01:ex-trajectory} asks you to sketch such a path yourself.}
\label{ch01:fig-trajectory}
\end{figure}

\begin{intuitionbox}
A single-channel signal is a bead sliding along a wire (one number, up and down). A
three-channel accelerometer is a bead flying around inside a room (three numbers, a 3D
trajectory). The richer the channel count, the more room the bead has to move -- and the more
ways a gesture can distinguish itself from every other gesture. Recognition gets easier with
more channels, up to a point.
\end{intuitionbox}

\begin{connectionbox}
Treating a multichannel signal as a vector-valued function is your first taste of a theme that
runs through the companion \emph{Linear Algebra From the Ground Up}: \emph{signals are vectors}.
In Chapter 3 we push this much further and treat an \emph{entire single-channel recording} -- all
its samples at once -- as one vector in a high-dimensional space. Inner products, orthogonality,
and projection, which you may know from that book as abstract operations, will turn out to be
exactly how the Fourier transform measures frequency. Hold the thought; it is the spine of the
whole subject.
\end{connectionbox}

\section{Analog versus digital}
\label{ch01:analog-digital}

Here is the deepest dividing line in the field, and the one most often misunderstood. It is
really \emph{two} independent distinctions wearing one name.

\subsection*{Continuous versus discrete in time}
A real microphone produces a voltage that exists at \emph{every} instant -- there is a value at
$t = 0.1$~s, at $t = 0.10001$~s, at every real number you can name. This is a
\textbf{continuous-time} signal: its domain is the real line, or an interval of it. We write
$x(t)$ with a parenthesis to flag it.

A computer cannot store a value for every one of the uncountably many instants. So we
\textbf{sample}: we look only at evenly spaced moments $t = 0, T_s, 2T_s, 3T_s, \dots$, where
$T_s$ is the \emph{sampling period}, and record the value at each. The result is a
\textbf{discrete-time} signal -- a sequence of numbers indexed by an integer $n$. We write
$x[n]$ with a square bracket to flag it; here $x[n]$ means "the value of the underlying signal at
time $n T_s$."

\begin{formalbox}
Notation we will keep for the rest of the book:
\begin{itemize}
\item Round brackets $x(t)$, $t \in \R$: \emph{continuous-time}. The domain is a continuum.
\item Square brackets $x[n]$, $n \in \Z$: \emph{discrete-time}. The domain is the integers.
\end{itemize}
The sampling rate is $f_s = 1/T_s$, in samples per second (often written Hz). A CD samples audio
at $f_s = 44100$~Hz; a typical wrist IMU might run at $f_s = 100$ to $1600$~Hz.
\end{formalbox}

\subsection*{Continuous versus discrete in value}
The second distinction is about the \emph{value}, not the time. A continuous-time voltage can sit
at \emph{any} real value: $0.7314\ldots$ volts, exactly. But to store a number a computer must
round it to one of finitely many levels -- this is \textbf{quantization}. A signal whose values
are restricted to a finite set of levels is \emph{quantized}.

Putting both together:
\begin{itemize}
\item \textbf{Analog} signal: continuous in \emph{both} time and value. The native output of
microphones, accelerometer sensing elements, electrodes.
\item \textbf{Digital} signal: discrete in \emph{both} time and value -- sampled \emph{and}
quantized. This is what lives in a file, a register, an array.
\end{itemize}
The device that crosses from analog to digital is the \textbf{analog-to-digital converter}
(ADC); the device that crosses back is the \textbf{digital-to-analog converter} (DAC).
Figure~\ref{ch01:fig-sampling} shows the two acts -- sampling and quantization -- happening to a
smooth analog signal.

\begin{figure}[H]
\centering
\begin{tikzpicture}
\begin{axis}[
  width=13cm, height=5.6cm,
  xlabel={time $t$}, ylabel={value},
  xmin=0, xmax=1, ymin=-1.35, ymax=1.55,
  axis lines=left,
  ytick={-1,-0.5,0,0.5,1},
  yticklabels={$-1$,,$0$,,$1$},
  legend style={at={(0.5,1.18)}, anchor=north, legend columns=3, draw=none, /tikz/every even column/.append style={column sep=0.4cm}},
  enlargelimits=false,
]
% analog curve
\addplot[sigblue, thick, domain=0:1, samples=200] {0.9*sin(deg(2*pi*1.5*x)) };
\addlegendentry{analog $x(t)$}
% sampled (stems) -- quantized to 0.25 grid
\addplot[specred, only marks, mark=*, mark size=1.6pt, samples at={0,0.0667,...,1.0}]
  ({x}, {0.25*round(0.9*sin(deg(2*pi*1.5*x))/0.25)});
\addlegendentry{samples (quantized)}
% stems
\addplot[specred, ycomb, samples at={0,0.0667,...,1.0}]
  ({x}, {0.25*round(0.9*sin(deg(2*pi*1.5*x))/0.25)});
\addlegendentry{}
\end{axis}
\end{tikzpicture}
\caption{From analog to digital. The smooth blue curve is the continuous-time, continuous-value
analog signal. Sampling keeps it only at evenly spaced instants (the dots); quantization snaps
each sample to the nearest level on a fixed value grid (here multiples of $0.25$). The dots are
the digital signal $x[n]$.}
\label{ch01:fig-sampling}
\end{figure}

\begin{trapbox}
"Digital is just analog but on a computer" -- this is the single most common misconception, and
it hides two real losses. Sampling can throw away fast wiggles that fall between the sample
instants (\emph{aliasing}, Chapter 18). Quantization adds a small rounding error to every value
(\emph{quantization noise}, Chapter 19). Both are usually manageable, but they are real, and
pretending the digital signal is identical to the analog one will eventually bite you. We will
make peace with both losses rigorously later; for now, just know they exist.
\end{trapbox}

\begin{historybox}
The word "signal" predates electronics by centuries -- think of smoke signals or naval flags,
discrete symbols sent across distance. The \emph{continuous} electrical signal arrived with the
telegraph and especially the telephone (1870s onward), which had to carry the continuously
varying pressure of speech down a wire. The \emph{digital} signal as we know it crystallized in
the 1920s--1940s at Bell Labs, where Harry Nyquist and later Claude Shannon worked out precisely
how many samples per second you need to capture a continuous signal without loss. That result --
the sampling theorem -- is the bridge this whole book is built on, and we will earn it carefully
in Chapter 18.
\end{historybox}

\begin{connectionbox}
The continuous/discrete split is exactly the boundary between integration and summation that you
met in \emph{Calculus From the Ground Up}. Almost every formula in this book comes in two
versions: a continuous one with an integral $\int x(t)\,\dd t$ and a discrete one with a sum
$\sum_n x[n]$. They are the same idea -- "accumulate the signal" -- viewed through the two lenses
of the calculus book. When a discrete formula looks strange, recover the intuition by asking
what its continuous integral cousin is doing.
\end{connectionbox}

\section{The five central questions of DSP}
\label{ch01:questions}

Why is there a whole field -- and a whole book -- about these wiggling functions? Because once
you have a signal, there are exactly five things you ever want to do with it. Every technique
you will learn is an answer to one of these.

\begin{enumerate}
\item \textbf{Represent.} How do we store and describe a signal accurately and efficiently? How
many samples? How many bits per sample? Can we describe it by its frequencies instead of its
samples? (Sampling, quantization, Fourier representations.)

\item \textbf{Transform.} How do we re-express the signal in a form that reveals what we care
about? The Fourier transform is the giant here: it rewrites a signal-in-time as a
signal-in-frequency, turning "what does it look like moment to moment" into "what frequencies is
it made of." (Chapters 6--11, 15.)

\item \textbf{Filter.} How do we keep the parts we want and remove the parts we do not -- the
hum, the drift, the high-frequency hiss, the gravity baseline in an accelerometer? (LTI systems,
convolution, filter design, Chapters 12--17.)

\item \textbf{Analyze.} How do we measure properties of the signal -- its dominant frequency,
its energy, how it changes over time, how two channels relate? (Spectral estimation,
spectrograms, statistics of noise, Chapters 10, 20--21.)

\item \textbf{Recognize.} How do we decide \emph{what the signal means}? Is this a heartbeat or
an artifact? Is this gesture a flick or a twist? This is where signal processing hands off to
pattern recognition and machine learning. (Chapters 38--39.)
\end{enumerate}

These five are not a random list; they line up into a pipeline, and our controller walks the
whole length of it on every gesture. Figure~\ref{ch01:fig-pipeline} draws that pipeline so you
can carry the picture into the chapters that fill in each stage.

\begin{figure}[H]
\centering
\begin{tikzpicture}[
  node distance=4mm,
  stage/.style={draw, rounded corners, align=center, font=\footnotesize\bfseries,
    minimum height=10mm, text width=18mm, fill=blue!4},
  flow/.style={-{Stealth[length=2mm]}, thick},
]
\node[stage] (rep) {Represent};
\node[stage, right=of rep] (tra) {Transform};
\node[stage, right=of tra] (fil) {Filter};
\node[stage, right=of fil] (ana) {Analyze};
\node[stage, right=of ana, fill=green!6] (rec) {Recognize};
\draw[flow] (rep) -- (tra);
\draw[flow] (tra) -- (fil);
\draw[flow] (fil) -- (ana);
\draw[flow] (ana) -- (rec);
\node[font=\scriptsize, below=2mm of rep, text width=20mm, align=center] {sample the
IMU stream};
\node[font=\scriptsize, below=2mm of tra, text width=20mm, align=center] {window to
frequency};
\node[font=\scriptsize, below=2mm of fil, text width=20mm, align=center] {drop gravity,
jitter};
\node[font=\scriptsize, below=2mm of ana, text width=20mm, align=center] {burst energy,
peak freq.};
\node[font=\scriptsize, below=2mm of rec, text width=20mm, align=center] {emit MIDI for
the gesture};
\end{tikzpicture}
\caption{The five central questions as one pipeline. The first four stages progressively refine
the signal -- raw samples, a frequency picture, a cleaned version, a few measured numbers -- and
the last stage makes a decision. The bottom row shows what each stage does inside the wrist
controller.}
\label{ch01:fig-pipeline}
\end{figure}

\begin{happeningbox}
Notice the arc. The first four questions progressively \emph{refine} the signal -- from raw
samples, to a frequency picture, to a cleaned-up version, to a handful of measured numbers. The
fifth question makes a \emph{decision}. Our wrist controller walks this exact path on every
gesture: represent the IMU stream as samples; transform a window of it to see its frequency
content; filter out gravity and jitter; analyze for features like burst energy and dominant
frequency; and finally recognize which gesture occurred and emit a MIDI message. The book is
ordered, roughly, to march you down this pipeline.
\end{happeningbox}

\section{A first clear-eyed look at the wrist controller}
\label{ch01:limits}

It is tempting -- and you will see it claimed in marketing -- that a wristband can read your
individual finger movements from motion alone. It is worth being precise about this from page
one, because being clear about what a sensor can and cannot do is the difference between an
engineer and a salesperson.

\begin{goalbox}
\textbf{What a wrist-worn IMU can and cannot do.} An inertial measurement unit (IMU) on the
wrist measures the \emph{motion of the wrist} -- its linear acceleration and its rotation rate.
From that it can reliably detect \emph{coarse, whole-hand gestures}: a flick, a twist, a tap, a
shake, raising or lowering the forearm, a circular wave. These are discrete, distinguishable
motions of the limb, and they are plenty to drive an expressive instrument.

What a wrist IMU \emph{cannot} do is recover the continuous joint angles of your individual
fingers. Curling your index finger while holding your wrist still produces almost no signal at
the wrist; the IMU is, quite literally, in the wrong place to see it. Recovering finger-level
activity from a wristband requires a \emph{different physical channel}: electromyography (EMG) or
force myography (FMG), which sense the forearm muscles and tendons that drive the fingers and run
right under a wristband. Even those give muscle \emph{activation}, not clean joint angles;
continuous finger \emph{pose} realistically needs sensors on the fingers themselves (instrumented
gloves or rings). We will treat EMG and FMG with the same care in Chapter 37 and cite the real
research there. For now, set your expectations correctly: \emph{the IMU gives us gestures of the
hand as a whole; finger-level detail needs help.}
\end{goalbox}

This is not a disappointment -- it is a design constraint, and good engineering lives inside
constraints. A flick, a twist, a tap, and a hold, mapped to musical actions, already make a real
instrument. Knowing the limit up front keeps us from chasing impossible signal processing later.

\begin{buildbox}
You can start collecting the very signals this book is about with almost no hardware. A
$\$5$--$\$15$ IMU breakout board (for example an MPU-6050, which combines a three-axis
accelerometer and three-axis gyroscope) wired to an Arduino over I\textsuperscript{2}C will
stream six channels of motion. The sketch below reads the accelerometer and prints three numbers
per line at a fixed rate -- a multichannel, discrete-time signal you can capture and plot. Do not
worry about the details yet; just notice that the output is exactly the object of
Definition~\ref{ch01:signal-defn}: values versus time, sample by sample.
\end{buildbox}

\begin{lstlisting}[language=C, caption={Streaming a 3-axis accelerometer as a discrete-time, multichannel signal (Arduino + MPU-6050).}]
#include <Wire.h>
const int MPU = 0x68;          // I2C address of the IMU
const int FS_HZ = 100;         // sampling rate: 100 samples/sec
const unsigned long DT_US = 1000000UL / FS_HZ;

void setup() {
  Serial.begin(115200);
  Wire.begin();
  Wire.beginTransmission(MPU);
  Wire.write(0x6B); Wire.write(0); // wake the sensor
  Wire.endTransmission(true);
}

void loop() {
  static unsigned long next = 0;
  unsigned long now = micros();
  if (now < next) return;        // hold a steady sampling period T_s
  next = now + DT_US;

  Wire.beginTransmission(MPU);
  Wire.write(0x3B);              // first accel register
  Wire.endTransmission(false);
  Wire.requestFrom(MPU, 6, true);
  int16_t ax = Wire.read() << 8 | Wire.read();
  int16_t ay = Wire.read() << 8 | Wire.read();
  int16_t az = Wire.read() << 8 | Wire.read();

  // One row = one sample x[n] = (ax, ay, az): three channels at instant n.
  Serial.print(ax); Serial.print(',');
  Serial.print(ay); Serial.print(',');
  Serial.println(az);
}
\end{lstlisting}

\begin{labbox}
Run the sketch (or grab any logged IMU dataset) and plot the three columns against sample index.
Then \emph{flick your wrist} and look at the trace. You will see a burst: the curves leap, swing,
and settle, much like Figure~\ref{ch01:fig-gallery}(b). Try a slow twist and a sharp tap, and
compare their shapes. You are already doing the first job of the whole field -- looking at a
signal and asking what its shape means. Everything technical in this book is in service of
answering that question reliably and automatically.
\end{labbox}

\section{Energy, power, and why we will care}
\label{ch01:energy}

One last pair of ideas will recur so often that it belongs in this opening chapter: how
"large" a signal is. There are two natural measures, and which one applies depends on whether the
signal eventually dies away or runs forever.

The \textbf{energy} of a signal is the running total of its value squared:
\[
  E = \int_{-\infty}^{\infty} \abs{x(t)}^2 \,\dd t
  \qquad\text{(continuous)}, \qquad
  E = \sum_{n=-\infty}^{\infty} \abs{x[n]}^2
  \qquad\text{(discrete)}.
\]
Squaring does two jobs at once: it makes negative dips count as much as positive peaks (a signal
that swings down hard is not "small"), and it weights big excursions heavily. A transient like a
single wrist flick has finite energy -- it starts, happens, and stops.

A signal that never stops -- a pure tone running forever, or the steady hum of mains power -- has
\emph{infinite} energy, so energy is the wrong yardstick for it. For these we use \textbf{power},
the energy per unit time, i.e.\ the average of $\abs{x}^2$:
\[
  P = \lim_{T\to\infty} \frac{1}{T}\int_{-T/2}^{T/2} \abs{x(t)}^2 \,\dd t .
\]

\begin{intuitionbox}
Energy is "how much total commotion was there, start to finish" -- the right question for a burst
that begins and ends, like a gesture. Power is "how much commotion per second, on average" -- the
right question for something ongoing, like background vibration or a sustained tone. A flick has
energy; a tremor has power.
\end{intuitionbox}

\begin{connectionbox}
The expression $\sum \abs{x[n]}^2$ is the squared \emph{length} of the signal-as-vector from the
linear-algebra book, and $\int \abs{x(t)}^2 \dd t$ is its continuous analogue. So "energy" is
just "length squared," and the famous Parseval/Plancherel theorem of Chapter 11 will say
something beautifully geometric: a signal's energy is the same whether you measure it in time or
in frequency, because the Fourier transform is a rotation that preserves length. You are
collecting the pieces of that picture already.
\end{connectionbox}

\begin{aestheticbox}
Step back and admire the unification. Sound, images, heartbeats, and the motion of your hand are,
to a wildly different set of human senses, utterly unlike one another. Yet we have just put them
all in one frame: each is a function from a domain (time or space) to a value, carrying
information in its variation, with a length we call energy. That one frame is what lets a single
toolkit -- the Fourier transform, filtering, sampling -- serve audio engineers, radiologists,
cardiologists, and instrument-builders alike. Seeing the common skeleton under different flesh is
the first real act of signal-processing thinking.
\end{aestheticbox}

\section{Looking ahead}
\label{ch01:ahead}

We now have the vocabulary the rest of the book assumes. A signal is a function carrying
information in its variation (Section~\ref{ch01:defn}). Its domain may be 1D time, 2D space, or
more (Section~\ref{ch01:dimension}); its value may be a single number or a vector of channels
(Section~\ref{ch01:channels}); and it may be analog (continuous in time and value) or digital
(sampled and quantized), Section~\ref{ch01:analog-digital}. Everything we do answers one of five
questions -- represent, transform, filter, analyze, recognize (Section~\ref{ch01:questions}) --
and our concrete target throughout is a wrist gesture controller whose real limits we have
already stated (Section~\ref{ch01:limits}).

The very next chapter sharpens the continuous-versus-discrete distinction into precise
mathematical objects we can compute with. Then, in Chapter 3, we take the leap that powers the
entire subject: we stop thinking of a signal as a wiggling curve and start thinking of it as a
single \emph{vector} in a vast space -- the geometric idea from which frequency and the Fourier
transform will emerge as naturally as shadows from a light.

\section{Exercises}
\label{ch01:exercises}

\begin{exercise}
For each of the following, state (i) the dimension of the domain, (ii) the number of channels,
and (iii) whether it is most naturally modeled as analog or digital: a vinyl record groove; a
stereo MP3 file; a black-and-white scanned document; the raw output of a three-axis gyroscope; a
color video clip; the reading on a digital kitchen scale once per second.
\end{exercise}

\begin{exercise}
Explain, in your own words and without formulas, why a constant signal $x(t)=c$ carries no
information, while $x(t)=\sin(2\pi f t)$ does. What is the smallest change you could make to the
constant signal so that it begins to carry information?
\end{exercise}

\begin{exercise}
A wrist IMU is sampled at $f_s = 200$~Hz for $3$ seconds while you perform one flick. (a) How
many samples are recorded per channel? (b) If it has $6$ channels (accel + gyro) stored as
$16$-bit integers, how many bytes does the recording occupy? (c) What is the sampling period
$T_s$ in milliseconds?
\end{exercise}

\begin{exercise}
Sketch (by hand or with code) a discrete-time signal $x[n]$ for $n=0,1,\dots,15$ that you would
expect from a single sharp tap on a table, recorded by one accelerometer axis. Describe in words
the shape you drew and why.
\end{exercise}

\begin{exercise}
Quantization snaps each sample to the nearest level on a value grid. If an ADC has $8$ bits and
covers the range $[-2g, +2g]$ (where $g$ is the acceleration of gravity), how many levels are
there, and what is the spacing between adjacent levels (the resolution) in units of $g$? Repeat
for a $12$-bit ADC over the same range and comment on the improvement.
\end{exercise}

\begin{exercise}
Classify each task below under exactly one of the five central questions (represent, transform,
filter, analyze, recognize): (a) removing the constant gravity component from an accelerometer
axis; (b) deciding whether a gesture was a flick or a twist; (c) choosing how many samples per
second to store; (d) computing the dominant frequency of a tremor; (e) re-expressing a chord as
the set of frequencies it contains.
\end{exercise}

\begin{exercise}
Consider the finite discrete signal $x[n] = \{2, -1, 3, 0, -2\}$ for $n=0,\dots,4$ and zero
elsewhere. Compute its energy $E=\sum_n \abs{x[n]}^2$. Then form $y[n]=2x[n]$ and compute its
energy. By what factor did the energy change, and why?
\end{exercise}

\begin{exercise}
A pure tone $x(t)=A\sin(2\pi f t)$ runs forever. Argue that its energy is infinite but its power
is finite. (You do not need to evaluate the integral exactly; reason about the average of
$\sin^2$ over a full cycle, which is $1/2$.) What is the power in terms of $A$?
\end{exercise}

\begin{exercise}
A marketing brochure claims a wrist band "tracks every finger in 3D using its built-in motion
sensor." Using the discussion in Section~\ref{ch01:limits}, write two or three sentences
explaining what part of this claim is physically plausible and what part is not, and name the
kind of sensor that would actually be needed for the implausible part.
\end{exercise}

\begin{exercise}
\label{ch01:ex-trajectory}
A multichannel signal can be viewed as a moving point in $\R^C$. For a two-axis sensor
($C=2$) recording a smooth circular wave of the hand, sketch the path traced by the point
$(x_1(t),x_2(t))$ in the plane (not versus time -- the two axes plotted against \emph{each
other}). What shape do you expect, and what would a faster wave change about it?
\end{exercise}

\begin{exercise}
(Open-ended.) Record or download a short IMU trace of yourself making three different gestures.
Plot all channels. Without doing any math, list three \emph{visual} features that distinguish the
gestures (for example: which axis spikes first, how long the burst lasts, whether it oscillates).
These hand-found features are exactly what an automatic recognizer will later learn -- you are
sketching its job by eye.
\end{exercise}
